% OC Core 1.3.3 -- module architecture and cross-level structures

\section{Module Architecture and Cross-Level Structures}
\label{sec:modules_master}

This section provides a high-level overview of the extended structural modules
of the Ontology of Continua. Each module is developed in the technical atlas;
here we summarise their roles and relations so the main spine has a complete
bridge between formal doctrine and the Core~1.3.3 closure chapters.

The extended modules are:

\begin{itemize}
\item \textbf{K-levels master} (the K-level overview chapter):
  global overview of declared continuum classes \(K_0\)–\(K_{12}\), their state spaces,
  axes, thresholds and core processes.

\item \textbf{Embedding spaces / M-spaces} (the M-space overview chapter):
  definition of meta-spaces \(M_x\) that constrain and parameterise the
  evolution of continua \(K_x\).

\item \textbf{Cross-level / cross-K structures} (the cross-level structure chapter):
  operators, mappings and landscape structures that link different levels
  of the hierarchy and generate phase transitions between them.

\item \textbf{Processes} (the process taxonomy chapter):
  unified description of evolution operators \(E = (F,G,H,Q,R,S,U)\) across
  levels, including tension dynamics, thresholds, flows and cycles.

\item \textbf{Cycles} (the cycle taxonomy chapter):
  structural and dynamical cycles \(C_j\) as carriers of continuity,
  memory and temporal organisation.

\item \textbf{Predictions} (the bounded prediction-route chapter):
  level-specific and cross-level predictions, empirical signatures and
  qualitative regimes implied by the formalism.

\item \textbf{Experiments} (the experiment-design chapter):
  conceptual and concrete experimental designs to probe continua,
  transitions and threshold landscapes.

\item \textbf{Falsifiability structure} (the falsifiability chapter):
  explicit conditions under which the Ontology of Continua can be
  refuted, constrained or forced to restructure its core.

\item \textbf{Jets / local evolution structure} (the local-evolution chapter):
  local expansions of the evolution operator and short-time dynamics
  around given configurations of \(K_x\) and \(M_x\).
\end{itemize}

Together, these modules refine and operationalise the core formalism across
levels and applications, while maintaining a single coherent hierarchy of
continua and meta-spaces.
